\chapter{Osilator Harmonik Kuantum}\label{ch:b3:harmonic-oscillator}

Hampir tak ada apa pun di alam yang benar-benar osilator harmonik, dan
hampir segala sesuatu merupakan osilator harmonik secara hampiran:
sistem apa pun yang disenggol dari kesetimbangan mantapnya --- ikatan
sebuah molekul, sebuah atom dalam kristal, sebuah jembatan, sebuah ragam
medan elektromagnetik --- merasakan gaya pemulih yang sebanding dengan
simpangannya, karena setiap potensial yang mulus berupa parabola di
dasar sumurnya. Siapa pun yang menyelesaikan osilator kuantum sekali
karena itu menyelesaikan getaran kecil seluruh dunia. Bab ini
menyelesaikannya dalam gaya aljabar yang telah menjadi ciri khas
mekanika kuantum: dua \omterm{def:b3:harmonic-oscillator:ladder}{operator tangga} memanjat dan menuruni tangga aras
$\big(n +
\tfrac12\big)\hbar\omega$ yang serba rata, dengan setengah anak tangga
di dasarnya --- yakni \omterm{thm:b3:harmonic-oscillator:spectrum}{energi titik nol} --- yang merupakan sebuah teorema,
bukan pilihan. Akibatnya membentang dari mengapa helium tak pernah
membeku sampai bagaimana molekul karbon dioksida, yang berdenting pada
$\hbar\omega$-nya sendiri, mencegat kalor Bumi yang keluar.

\section{Tangganya}

\begin{definition}[Operator tangga]\label{def:b3:harmonic-oscillator:ladder}
\index{operator tangga}\index{operator pemusnah}\index{operator pencipta}
Untuk $\hat H = \hat p^2/2m + \tfrac12 m\omega^2\hat x^2$,
perkenalkanlah pasangan tak berdimensi yang tak Hermitian
\[
\hat a = \sqrt{\frac{m\omega}{2\hbar}}\Big(\hat x +
\frac{\iu\hat p}{m\omega}\Big) , \qquad
\hat a^\dagger = \sqrt{\frac{m\omega}{2\hbar}}\Big(\hat x -
\frac{\iu\hat p}{m\omega}\Big) .
\]
Dari $[\hat x, \hat p] = \iu\hbar$:
\[
[\hat a, \hat a^\dagger] = 1 , \qquad
\hat H = \hbar\omega\Big(\hat N + \tfrac12\Big) , \quad
\hat N = \hat a^\dagger\hat a .
\]
$\hat N$ bersifat Hermitian; \omterm{def:b3:quantum-formalism:observable}{nilai eigennya} akan mencacah \emph{kuantum},
dan $\hat a$, $\hat a^\dagger$ --- yakni operator \emph{pemusnah} dan
\emph{pencipta} --- akan membuang dan menambahkan satu.
\end{definition}

\begin{theorem}[Spektrumnya, hanya lewat aljabar]\label{thm:b3:harmonic-oscillator:spectrum}
\index{energi titik nol}
\omterm{def:b3:quantum-formalism:observable}{Nilai eigen} $\hat H$ tepat sama dengan
\[
E_n = \Big(n + \tfrac12\Big)\hbar\omega , \qquad n = 0, 1, 2, \dots
\]
--- yakni tangga tak hingga beranak tangga sama $\hbar\omega$ di atas
aras dasar yang \emph{bukan} nol: yaitu \emph{\omterm{thm:b3:harmonic-oscillator:spectrum}{energi titik nol}}
$\tfrac12\hbar\omega$. Keadaan eigen ternormalkan $\ket n$ terhubung
oleh
\[
\hat a\ket n = \sqrt n\,\ket{n - 1} , \qquad
\hat a^\dagger\ket n = \sqrt{n + 1}\,\ket{n + 1} ,
\qquad
\ket n = \frac{(\hat a^\dagger)^n}{\sqrt{n!}}\,\ket0 .
\]
\end{theorem}

\begin{proof}
Dari \omterm{def:b3:quantum-formalism:commutator}{komutatornya}, $\hat N\hat a = \hat a(\hat N - 1)$: jika
$\ket\nu$ mempunyai \omterm{def:b3:quantum-formalism:observable}{nilai eigen} $\hat N$ sebesar $\nu$, maka $\hat a\ket\nu$
adalah vektor eigen dengan $\nu - 1$ (atau vektor nol). Tiap penurunan
hanya diizinkan selama $\nu \ge 0$, karena $\nu = \braket{\nu}{\hat
N\nu} = \|\hat a\ket\nu\|^2 \ge 0$; penurunannya karena itu harus
berhenti, dan ia hanya berhenti pada keadaan dengan $\hat a\ket{\nu_0}
= 0$, sehingga $\nu_0 = 0$. Jadi $\nu$ menjangkau bilangan bulat tak
negatif. Penormalannya menyusul dari $\|\hat a\ket n\|^2 = n$
dan $\|\hat a^\dagger\ket n\|^2 = n + 1$.
\end{proof}

\begin{omfigure}
\begin{tikzpicture}[font=\small, >=Stealth]
  \draw[->] (0,-0.3) -- (0,5.4);
  \node[anchor=south] at (0,5.4) {$E$};
  \foreach \n in {0,1,2,3,4}
    {\draw[very thick, omDef] (0.5,{0.55+\n*1.05}) -- (3.3,{0.55+\n*1.05});
     \node[anchor=west] at (3.45,{0.55+\n*1.05})
       {$\ket{\n}$\; $E = (\n + \tfrac12)\hbar\omega$};}
  \draw[->, omThm, very thick] (1.35,0.62) -- (1.35,1.53);
  \node[omThm, anchor=east] at (1.28,1.08) {$\hat a^\dagger$};
  \draw[->, omExo, very thick] (2.45,1.53) -- (2.45,0.62);
  \node[omExo, anchor=west] at (2.52,1.08) {$\hat a$};
  \draw[<->] (0.15,0) -- (0.15,0.55);
  \node[anchor=west, font=\footnotesize] at (0.28,0.22) {$\tfrac12\hbar\omega$};
  \draw[dashed] (0,0) -- (3.3,0);
\end{tikzpicture}

{\small Tangga osilatornya: anak tangga yang sama besar $\hbar\omega$,
yang dipanjat $\hat a^\dagger$ dan dituruni $\hat a$, berdiri di atas
lantai setengah anak tangga di atas \omterm{def:b3:relativistic-dynamics:momentum-energy}{energi diam} klasiknya.}
\end{omfigure}

\begin{proposition}[Keadaannya dalam ruang]\label{prop:b3:harmonic-oscillator:wavefunctions}
Keadaan dasarnya adalah fungsi Gauss
\[
\varphi_0(x) = \Big(\frac{m\omega}{\pi\hbar}\Big)^{1/4}
\eu^{-m\omega x^2/2\hbar} ,
\]
yang diperoleh dengan menyelesaikan persamaan \emph{orde satu}
$\hat a\varphi_0
= 0$; ia menjenuhkan ketaksamaan Heisenberg, $\Delta x\,\Delta p =
\hbar/2$, dengan $\Delta x = \sqrt{\hbar/2m\omega}$. Mengenakan
$\hat a^\dagger$ berulang kali membangkitkan $\varphi_n$: sebuah
polinomial berderajat $n$ (dengan $n$ simpul) dikalikan fungsi Gauss
yang sama. Untuk $n$ besar, $|\varphi_n|^2$ berosilasi cepat di sekitar
sebaran waktu tinggal klasiknya, dan terbesar di dekat titik baliknya
--- yakni asas korespondensi dalam sebuah gambar.
\end{proposition}

\begin{proof}[Bukti sebagian]
$\hat a\varphi_0 = 0$ berbunyi $\varphi_0' = -(m\omega/\hbar)x
\varphi_0$: yakni fungsi Gauss, yang ternormalkan. Penjenuhannya: fungsi
Gauss adalah kasus kesamaan
\cref{thm:b3:quantum-formalism:uncertainty}. Struktur polinomialnya
menyusul karena $\hat a^\dagger$ berorde satu dalam $x$ dan $\dd/\dd x$;
pernyataan untuk $n$ besarnya diperiksa dalam
\cref{exo:b3:harmonic-oscillator:12}.
\end{proof}

\begin{omfigure}
\begin{tikzpicture}
  \begin{axis}[omaxis, axis x line=bottom, axis y line=left,
    width=6.6cm, height=5.2cm, xmin=-4, xmax=4, ymin=0, ymax=3.6,
    xlabel={$x$ (dalam satuan $\sqrt{\hbar/m\omega}$)}, ylabel={},
    xlabel style={at={(axis description cs:0.5,-0.1)}, anchor=north},
    xtick={-2,0,2}, ytick=\empty, clip=false]
    \addplot[omDef, thick, domain=-4:4, samples=90]
      {0.56*exp(-x*x)+0.25};
    \addplot[omThm, thick, domain=-4:4, samples=90]
      {1.13*x*x*exp(-x*x)+1.3};
    \addplot[omExo, thick, domain=-4:4, samples=120]
      {0.28*(2*x*x-1)^2*exp(-x*x)+2.35};
    \node[omDef, font=\scriptsize, anchor=west] at (axis cs:2.4,0.5) {$|\varphi_0|^2$};
    \node[omThm, font=\scriptsize, anchor=west] at (axis cs:2.4,1.55) {$|\varphi_1|^2$};
    \node[omExo, font=\scriptsize, anchor=west] at (axis cs:2.4,2.6) {$|\varphi_2|^2$};
  \end{axis}
  \begin{axis}[omaxis, axis x line=bottom, axis y line=left,
    width=6.6cm, height=5.2cm, xmin=-5.2, xmax=5.2, ymin=0, ymax=0.45,
    xshift=7.1cm,
    xlabel={$x$}, ylabel={},
    xlabel style={at={(axis description cs:0.5,-0.1)}, anchor=north},
    xtick=\empty, ytick=\empty, clip=false]
    \addplot[omDef, domain=-4.6:4.6, samples=300]
      {(64*x^6-480*x^4+720*x^2-120)^2*exp(-x*x)/81680};
    \addplot[omThm, very thick, dashed, domain=-3.54:3.54, samples=120]
      {0.3183/sqrt(max(0.35,13-x*x))};
    \node[omThm, font=\scriptsize, anchor=south] at (axis cs:0,0.36) {waktu tinggal klasik};
    \node[omDef, font=\scriptsize, anchor=west] at (axis cs:1.9,0.42) {$|\varphi_6|^2$};
  \end{axis}
\end{tikzpicture}

{\small Kiri: rapat peluang terendahnya (yang digeser tegak):
sebuah fungsi Gauss tanpa simpul, lalu $n$ simpul untuk $\varphi_n$.
Kanan: pada $n$ besar rapat kuantumnya berosilasi di sekitar sebaran
klasiknya, yang menumpuk di titik baliknya, tempat massa yang berosilasi
itu berlama-lama.}
\end{omfigure}

\section{Energi titik nol itu nyata}

\begin{remark}[Mengapa lantainya tak dapat lebih rendah]\label{rem:b3:harmonic-oscillator:zeropoint}
Keadaan berenergi nol akan menuntut $\langle\hat p^2\rangle =
\langle\hat x^2\rangle = 0$: benar-benar diam \emph{dan} benar-benar
terpusat, yang dilarang $\Delta x\,\Delta p \ge \hbar/2$. Keadaan
dasarnya adalah kompromi terbaik yang diizinkan ketaksamaan itu, dan
$\tfrac12\hbar\omega$ adalah sewanya. Ia bukan tetapan pembukuan: gerak
titik nol mengaburkan pola difraksi sinar-X pada nol mutlak; ia membuat
isotop yang lebih ringan mempunyai ikatan efektif yang lebih lemah
(H$_2$ dan D$_2$ terurai pada energi yang berbeda secara terukur dari
sumur elektronik yang sama); dan pada helium ia begitu ganas --- atomnya
ringan, tarikannya lemah --- sehingga zat cairnya \emph{tak pernah}
membeku di bawah tekanan uapnya sendiri: satu-satunya unsur yang masih
cair pada nol mutlak, dan hanya memadat dengan bantuan tekanan 25
atmosfer.
\end{remark}

\begin{example}[Skala $\hbar\omega$]\label{ex:b3:harmonic-oscillator:scales}
Ikatan karbon monoksida ($\omega = \qty{4.1e14}{rad/s}$): $\hbar\omega =
\qty{0.27}{eV}$, sepuluh kali $k_{\text{B}}T$ suhu ruang ---
getaran molekul membeku dalam udara sehari-hari, dan itulah sebabnya
kapasitas kalor gas dwiatom membingungkan abad kesembilan belas. Sebuah
bandul ($\omega = \qty{5}{rad/s}$): $\hbar\omega = \qty{3e-15}{eV}$,
sama sekali di luar jangkauan resolusi --- yakni limit korespondensinya.
Cermin LIGO seberat \qty{40}{kg}, yang digantung agar berayun pada
\qty{100}{Hz}, mempunyai sebagai ragam osilatornya amplitudo titik nol $\Delta x =
\sqrt{\hbar/2m\omega} \approx \qty{5e-20}{m}$ --- dan observatoriumnya
secara rutin memilah perpindahan \emph{tepat} pada lantai kuantum ini:
gerak titik nol benda seberat empat puluh kilogram kini menjadi kekangan
rekayasa (\cref{exo:b3:harmonic-oscillator:3}).
\end{example}

\section{Keadaan koheren: wajah klasik osilator kuantum}

\begin{definition}[Keadaan koheren]\label{def:b3:harmonic-oscillator:coherent}
\index{keadaan koheren}
Yang disebut \emph{keadaan koheren} $\ket\alpha$ adalah vektor eigen
\omterm{def:b3:harmonic-oscillator:ladder}{operator pemusnah}, $\hat a\ket\alpha = \alpha\ket\alpha$, dengan
$\alpha$ sembarang bilangan kompleks. Dijabarkan pada tangganya,
\[
\ket\alpha = \eu^{-|\alpha|^2/2}\sum_{n=0}^\infty
\frac{\alpha^n}{\sqrt{n!}}\,\ket n :
\]
cacah kuantumnya tersebar Poisson dengan rata-rata $\langle\hat
N\rangle = |\alpha|^2$ dan rentang $\Delta N = |\alpha|$.
\end{definition}

\begin{proposition}[Mengapa laser bersifat klasik]\label{prop:b3:harmonic-oscillator:coherent-motion}
\omterm{def:b3:harmonic-oscillator:coherent}{Keadaan koheren} adalah fungsi Gauss keadaan dasar yang tergeser dalam
\omterm{def:b3:hamiltonian-mechanics:hamiltonian}{ruang fase}; di bawah perkembangan osilatornya ia \emph{tetap} koheren,
$\alpha(t) = \alpha\,\eu^{-\iu\omega t}$: pusatnya menjalankan tepat
gerak klasiknya, $\langle\hat x\rangle(t) =
x_0\cos\omega t + \cdots$, sedangkan lebarnya tetap minimum
$\Delta x\,\Delta p = \hbar/2$ selamanya --- yakni paket gelombang yang
tak pernah melebar. \omterm{def:b3:harmonic-oscillator:coherent}{Keadaan koheren} adalah cara osilator kuantum
menyamar sebagai osilator klasik; cahaya sebuah laser adalah \omterm{def:b3:harmonic-oscillator:coherent}{keadaan
koheren} sebuah ragam medan, dengan cacah fotonnya Poisson (yaitu
\emph{derau tembakan} tiap fotodetektor), dan medannya berosilasi
seperti kata Maxwell.
\end{proposition}

\begin{proof}[Bukti sebagian]
Penjabarannya menyusul dengan menuliskan $\ket\alpha = \sum c_n\ket n$
dalam $\hat a\ket\alpha = \alpha\ket\alpha$: $c_n = \alpha c_{n-1}/\sqrt
n$. Perkembangannya: tiap $\ket n$ memungut $\eu^{-\iu(n + 1/2)\omega t}$,
yang terjumlahkan kembali menjadi \omterm{def:b3:harmonic-oscillator:coherent}{keadaan koheren} $\alpha\eu^{-\iu\omega t}$
(kecuali fase globalnya). $\langle\hat x\rangle \propto
\operatorname{Re}\alpha(t)$ dari $\hat x \propto \hat a +
\hat a^\dagger$. Bahwa keadaannya adalah fungsi Gauss yang tergeser
diterima tanpa bukti di sini.
\end{proof}

\begin{omfigure}
\begin{tikzpicture}[font=\small, >=Stealth]
  \draw[->] (-3.1,0) -- (3.4,0) node[right] {$\langle\hat x\rangle$};
  \draw[->] (0,-3.0) -- (0,3.2) node[above] {$\langle\hat p\rangle$};
  \draw[gray, dashed] (0,0) circle (2.2);
  \fill[omDef!25, draw=omDef, thick] (2.2,0) circle (0.42);
  \fill[omDef!25, draw=omDef, thick] (50:2.2) circle (0.42);
  \fill[omDef!25, draw=omDef, thick] (100:2.2) circle (0.42);
  \draw[->, thick] (2.53,0.9) arc (20:40:2.6);
  \node[anchor=west] at (2.62,1.35) {$\omega$};
  \fill[omThm!30, draw=omThm, thick] (0,0) circle (0.42);
  \node[omThm, font=\footnotesize, anchor=east] at (-0.3,-0.68) {$\ket0$};
  \node[omDef, font=\footnotesize, anchor=west] at (2.3,-0.62) {$\ket\alpha$};
  \node[font=\footnotesize, anchor=west] at (-3.05,2.75)
    {gumpalan minimum $\Delta x\,\Delta p = \hbar/2$, beredar secara klasik};
\end{tikzpicture}

{\small Potret \omterm{def:b3:hamiltonian-mechanics:hamiltonian}{ruang fase} (bandingkanlah
\cref{ch:b3:hamiltonian-mechanics}): keadaan dasarnya adalah gumpalan
ketaktentuan minimum di titik asal; \omterm{def:b3:harmonic-oscillator:coherent}{keadaan koheren} adalah gumpalan yang
sama yang tergeser, beredar pada $\omega$ tanpa berubah bentuk --- yakni
tiruan terbaik mekanika kuantum atas osilasi klasik.}
\end{omfigure}

\begin{method}[Aljabar osilator]\label{met:b3:harmonic-oscillator:method}
(1) Nyatakanlah apa pun yang ditanyakan dalam $\hat a$, $\hat a^\dagger$: $\hat x
= \sqrt{\hbar/2m\omega}\,(\hat a + \hat a^\dagger)$, $\hat p =
\iu\sqrt{m\hbar\omega/2}\,(\hat a^\dagger - \hat a)$. (2) Geserlah
$\hat a$ ke kanan dengan $[\hat a, \hat a^\dagger] = 1$; pakailah
$\hat a\ket0 = 0$. (3) Unsur matriksnya: $\hat x$ hanya menghubungkan
anak tangga bertetangga --- yakni kaidah pemilihan $\Delta n = \pm1$
dalam spektroskopi getaran. (4) Sembarang \omterm{def:b3:hamiltonian-mechanics:hamiltonian}{Hamiltonian} kuadratik
(osilator terkopel, rangkaian, ragam medan) terdiagonalkan menjadi
tangga yang saling bebas: carilah \omterm{prop:b3:lagrangian-mechanics:modes}{ragam normalnya} lebih dahulu, lalu
kuantumkan masing-masingnya. (5) Bilangan lebih dahulu:
$\hbar\omega$ berbanding $k_{\text{B}}T$ memutuskan apakah sistemnya
kuantum atau klasik sebelum aljabar apa pun dikerjakan.
\end{method}

\section{Latihan}

\begin{exercise}[$\star$]\label{exo:b3:harmonic-oscillator:1}
(a) Periksalah $[\hat a, \hat a^\dagger] = 1$ dari $[\hat x, \hat p] =
\iu\hbar$. (b) Periksalah $\hat H = \hbar\omega(\hat a^\dagger\hat a +
\tfrac12)$ lewat penjabaran langsung. (c) Baliklah untuk menyatakan
$\hat x$ dan $\hat p$. (d) Tunjukkanlah bahwa $\hat N$ Hermitian lalu
jelaskanlah mengapa $\hat a$ sendiri tak dapat menjadi \omterm{def:b3:quantum-formalism:observable}{besaran teramati}.
\end{exercise}

\begin{exercise}[$\star$]\label{exo:b3:harmonic-oscillator:2}
Pada keadaan eigen $\ket n$: (a) tunjukkanlah $\langle\hat x\rangle =
\langle\hat p\rangle = 0$; (b) hitunglah $\langle\hat x^2\rangle$ dan
$\langle\hat p^2\rangle$; (c) simpulkanlah $\Delta x\,\Delta p = (n +
\tfrac12)\hbar$; (d) periksalah nisbah virialnya $\langle E_k\rangle =
\langle E_p\rangle$.
\end{exercise}

\begin{exercise}[$\star$]\label{exo:b3:harmonic-oscillator:3}
Amplitudo titik nol $\Delta x = \sqrt{\hbar/2m\omega}$: hitunglah
untuk (a) molekul CO ($\mu = \qty{1.14e-26}{kg}$, $\omega =
\qty{4.1e14}{rad/s}$), dibandingkan dengan panjang ikatannya
\qty{0.11}{nm}; (b) sebuah atom hidrogen dalam zat padat ($m =
\qty{1.7e-27}{kg}$, $\hbar\omega = \qty{0.1}{eV}$); (c) sebuah cermin
LIGO ($m = \qty{40}{kg}$, $f = \qty{100}{Hz}$); (d) urutkanlah
ketiganya sebagai bagian dari ukuran sistemnya lalu berilah ulasan
tentang siapa yang ``kuantum''.
\end{exercise}

\begin{exercise}[$\star$]\label{exo:b3:harmonic-oscillator:4}
(a) Dengan memakai hubungan tangganya, hitunglah $\bra m\hat x\ket n$
lalu tunjukkanlah bahwa ia lenyap kecuali bila $m = n \pm 1$. (b)
Simpulkanlah kaidah pemilihan getarannya $\Delta n = \pm1$ bagi serapan
cahaya (koplingnya $\propto\hat x$). (c) Mengapa molekul beratom tak
sejenis (CO) menyerap cahaya inframerah sedangkan N$_2$ tidak? (d)
Molekul sungguhan menunjukkan garis ``nada atas'' yang lemah pada
$\approx 2\hbar\omega$: apa yang disingkapkannya tentang potensialnya?
\end{exercise}

\begin{exercise}[$\star\star$]\label{exo:b3:harmonic-oscillator:5}
(a) Selesaikanlah $\hat a\varphi_0 = 0$ sebagai persamaan diferensial
lalu normalkanlah. (b) Bangkitkanlah $\varphi_1$ dan $\varphi_2$ dengan
mengenakan $\hat a^\dagger$. (c) Periksalah $\varphi_1 \perp \varphi_0$
hanya lewat paritasnya. (d) Sketsalah ketiga rapatnya lalu periksalah
cacah simpulnya.
\end{exercise}

\begin{exercise}[$\star\star$]\label{exo:b3:harmonic-oscillator:6}
Pratayang Boltzmann. Sekumpulan osilator yang serupa pada suhu $T$
menempati aras $n$ dengan peluang $\propto
\eu^{-E_n/k_{\text{B}}T}$ (jilid Tahun 2, faktor Boltzmann). (a)
Tunjukkanlah bahwa bilangan kuantum rata-ratanya $\langle n\rangle =
1/(\eu^{\hbar\omega/k_{\text{B}}T} - 1)$. (b) Hitunglah untuk getaran
CO pada \qty{300}{K} dan pada \qty{2000}{K}. (c) Tunjukkanlah bahwa
energi rata-ratanya menuju $k_{\text{B}}T$ pada suhu tinggi (ekipartisi
diperoleh kembali) dan menuju $\tfrac12\hbar\omega$ pada suhu rendah.
(d) Pada suhu berapa getaran \qty{15}{\micro m} (lenturan karbon
dioksida) mempunyai $\langle n\rangle = 0.1$?
\end{exercise}

\begin{exercise}[$\star\star$]\label{exo:b3:harmonic-oscillator:7}
Statistik \omterm{def:b3:harmonic-oscillator:coherent}{keadaan koheren}. (a) Dari penjabaran $\ket\alpha$,
tunjukkanlah bahwa $\mathcal P(n)$ bersebaran Poisson dengan rata-rata
$|\alpha|^2$. (b) Tunjukkanlah $\langle\hat N\rangle = |\alpha|^2$ dan
$\Delta N = |\alpha|$. (c) Sebuah laser \qty{1}{mW} pada \qty{633}{nm}:
berapa foton tiap sekon, dan berapa fluktuasi nisbinya
$\Delta N/\langle N\rangle$ dalam satu sekon? (d) Derau tembakan:
tunjukkanlah bahwa nisbah derau terhadap isyarat arus fotonya turun
sebagai $1/\sqrt{\langle N\rangle}$ --- mengapa berkas yang terang
tampak mulus.
\end{exercise}

\begin{exercise}[$\star\star$]\label{exo:b3:harmonic-oscillator:8}
Perkembangan \omterm{def:b3:harmonic-oscillator:coherent}{keadaan koheren}. (a) Kenakanlah fase perkembangannya pada
penjabarannya lalu tunjukkanlah $\ket{\alpha(t)}$ dengan $\alpha(t) =
\alpha\eu^{-\iu\omega t}$ (kecuali fase globalnya). (b) Simpulkanlah
$\langle\hat x\rangle(t)$ dan $\langle\hat p\rangle(t)$ lalu
bandingkanlah dengan penyelesaian klasiknya. (c) Mengapa keadaan eigen
energi, walaupun stasioner, \emph{tidak} memerikan bandul yang berayun
--- ciri \omterm{def:b3:harmonic-oscillator:coherent}{keadaan koheren} yang manakah yang memerikannya? (d)
Perkirakanlah $|\alpha|$ untuk bandul sungguhan (\qty{10}{g},
beramplitudo \qty{10}{cm}, \qty{1}{Hz}) lalu berilah ulasan.
\end{exercise}

\begin{exercise}[$\star\star$]\label{exo:b3:harmonic-oscillator:9}
Rangkaian LC kuantum. Sebuah lingkar adiwarsa dengan $L =
\qty{10}{nH}$ dan $C = \qty{0.4}{pF}$ mengosilasikan muatan dan fluks
seperti $x$ dan $p$. (a) Berapa frekuensi resonansinya (jilid Tahun 1)
dan berapa kuantum $\hbar\omega$-nya dalam $\unit{\micro eV}$? (b) Di
bawah suhu berapa $k_{\text{B}}T \ll \hbar\omega$, sehingga rangkaiannya
berada dalam keadaan dasarnya? (c) Mengapa kubit adiwarsa dijalankan
dalam lemari pendingin pengenceran pada \qty{20}{mK}? (d) Fluktuasi
tegangan titik nolnya $\Delta V = \Delta q/C$ dengan $\Delta q =
\sqrt{\hbar\omega C/2}$: hitunglah.
\end{exercise}

\begin{exercise}[$\star\star\star$]\label{exo:b3:harmonic-oscillator:10}
Van der Waals dari gerak titik nol. Modelkanlah dua atom netral
berjarak $R$ sebagai dua osilator dipol yang serupa (bermuatan $e$,
bermassa $m$, berfrekuensi $\omega_0$) yang simpangannya terkopel oleh
energi dipol $\lambda\,x_1x_2$ dengan $\lambda = e^2/2\pi\varepsilon_0R^3$.
(a) Tunjukkanlah bahwa koordinat normalnya $(x_1 \pm x_2)/\sqrt2$
berosilasi pada
$\omega_\pm = \omega_0\sqrt{1 \pm \lambda/m\omega_0^2}$. (b) Energi
dasarnya adalah $\tfrac\hbar2(\omega_+ + \omega_-)$: jabarkanlah sampai
orde kedua dalam $\lambda$ lalu tunjukkanlah bahwa energi interaksinya
adalah
\[
\Delta E = -\frac{\hbar\lambda^2}{8m^2\omega_0^3}
\ \propto\ -\frac{1}{R^6} .
\]
(c) Mengapa tarikan ini bersifat semesta --- hadir bahkan antara atom
yang sama sekali tak berdipol tetap? (d) Hukum $1/R^6$ adalah gaya van
der Waals dalam koreksi gas nyata jilid Tahun 1: apa, secara mikroskopis,
yang ``berfluktuasi'' --- dan apa yang akan terjadi pada gaya ini dalam
dunia dengan $\hbar = 0$?
\end{exercise}

\begin{exercise}[$\star\star\star$]\label{exo:b3:harmonic-oscillator:11}
Pita samping sebuah ion terperangkap. Sebuah ion dalam perangkap
harmonik ($f =
\qty{1}{MHz}$) menyerap cahaya laser. Karena ionnya
bergerak, spektrum serapannya menunjukkan garis elektroniknya pada
$\nu_0$ yang diapit garis pada $\nu_0 \pm f$: yakni transisi yang
mengubah bilangan kuantum \emph{geraknya} sebesar $\mp1$ bersama
transisi elektroniknya. (a) Berapa
$\hbar\omega_{\text{perangkap}}$ dalam $\unit{neV}$, dan mengapa memilah
pita sampingnya menuntut garis yang sangat sempit? (b) Menggerakkan
garis $\nu_0 - f$ membuang satu kuantum gerak tiap daur: jelaskanlah
\emph{pendinginan pita samping} menuju keadaan dasarnya. (c) Setelah
$\langle n\rangle \approx
0$, pita samping bawahnya lenyap sama sekali --- mengapa
ketaksetangkupan itu merupakan bukti tercapainya keadaan dasar kuantum?
(d) Beginilah keadaan dasar gerak sebuah atom tunggal --- dan cermin
LIGO berskala kilogram, dengan cara lain --- disahkan: nyatakanlah
``suhu'' berapa yang telah dicapai ionnya untuk $f = \qty{1}{MHz}$ dan
$\langle n\rangle = 0.05$.
\end{exercise}

\begin{exercise}[$\star\star\star$]\label{exo:b3:harmonic-oscillator:12}
Limit klasiknya, secara kuantitatif. Osilator klasik beramplitudo
$A$ menghabiskan dalam $[x, x + \dd x]$ bagian sebesar $\dd
t/T = \dd x/\pi\sqrt{A^2 - x^2}$. (a) Turunkanlah sebaran waktu tinggal
ini. (b) Untuk keadaan kuantum $n$, ambillah $A_n$ dari $E_n =
\tfrac12 m\omega^2A_n^2$ lalu bandingkanlah sebaran klasiknya dengan
$|\varphi_n|^2$ yang dirata-ratakan secara setempat (yang diberikan): di
mana keduanya sesuai dan di mana keduanya harus berbeda? (c)
Tunjukkanlah bahwa jarak pecahan antara aras berdekatan, $\Delta E/E$,
lenyap sebagai $1/n$: energinya menjadi sinambung secara efektif. (d)
Untuk bandul pada \cref{exo:b3:harmonic-oscillator:8}(d),
perkirakanlah $n$ dan $\Delta
E/E$, lalu simpulkanlah argumen korespondensinya dalam satu kalimat.
\end{exercise}

\section{Soal: Molekul yang menghangatkan Bumi}

\begin{problem}[{Soal akhir pekan --- tangga kuantum karbon dioksida
dan efek rumah kaca}]\label{pb:b3:harmonic-oscillator:1}
Molekul karbon dioksida adalah rantai lurus O=C=O: beberapa osilator
terkuantumkan. Bahwa $\hbar\omega$ ragam lenturannya kebetulan berada di
tengah-tengah pijar termal Bumi yang keluar itulah sebabnya gas surih
ini --- empat molekul tiap sepuluh ribu --- mengemudikan iklim planet.
Data: bilangan gelombang lenturannya $\tilde\nu_2 = \qty{667}{cm^{-1}}$
(satuan para spektroskopis: $E = hc\tilde\nu$, $\qty{1}{cm^{-1}} =
\qty{1.24e-4}{eV}$); uluran taksetangkupnya $\tilde\nu_3 =
\qty{2349}{cm^{-1}}$; uluran setangkupnya $\tilde\nu_1 =
\qty{1388}{cm^{-1}}$; $k_{\text{B}}T$ pada $\qty{288}{K}$ adalah
\qty{24.8}{meV}; hukum Wien (jilid Tahun 2): $\lambda_{\max}T =
\qty{2898}{\micro m.K}$.

\medskip\noindent\textbf{Bagian I --- Ragam sebuah molekul lurus.}
\begin{enumerate}
  \item Tiga atom mempunyai sembilan \omterm{def:b3:lagrangian-mechanics:coordinates}{derajat kebebasan}: berapa yang
        merupakan translasi seluruh molekulnya, berapa rotasinya
        (molekulnya \emph{lurus}), dan berapa getaran yang tersisa?
  \item Perikanlah keempat getarannya: uluran setangkup, uluran
        taksetangkup, dan lenturan yang berdegenerasi ganda --- mengapa
        lenturannya muncul dua kali?
  \item Cahaya terkopel dengan \emph{dipol listrik} yang bergetar
        (\cref{exo:b3:harmonic-oscillator:4}). Ragam O=C=O yang manakah
        yang memodulasi momen dipolnya, dan yang mana yang bisu terhadap
        inframerah?
  \item Ubahlah ketiga bilangan gelombangnya menjadi panjang gelombang
        foton, lalu tempatkanlah: yang mana yang berada dalam inframerah
        termal?
  \item Mengapa kedua gas udara utama, N$_2$ dan O$_2$, praktis tidak
        menyerap inframerah sama sekali --- dan apa akibatnya bagi
        ketembusan atmosfernya?
  \item Uap air, yang bengkok dan berdipol, menyerap pada sebagian besar
        inframerah: pada ``jendela'' spektrum yang manakah lenturan
        karbon dioksida bekerja nyaris sendirian (bandingkanlah
        \qty{15}{\micro m} dengan pita kuat air di bawah
        \qty{8}{\micro m} dan di atas \qty{20}{\micro m})?
\end{enumerate}

\medskip\noindent\textbf{Bagian II --- Tangga kuantum lenturannya.}
\begin{enumerate}[resume]
  \item Hitunglah $\hbar\omega_2$ dalam meV, beserta tangganya $E_n$.
  \item Berapa bagian molekul yang menempati $n = 1$ pada
        \qty{288}{K} (terhadap $n = 0$, faktor Boltzmann)? Dan
        $n = 2$?
  \item Panjang gelombang foton yang manakah yang menggerakkan
        $n = 0 \to 1$? Mengapa panjang gelombang yang \emph{sama}
        menguasai pancarannya?
  \item Berilah alasan dari tangga harmoniknya bahwa satu panjang
        gelombang melayani seluruh tangganya ($n \to n + 1$ untuk tiap
        $n$): sifat jarak aras yang manakah yang bekerja?
  \item Molekulnya juga berputar, sehingga menambahkan struktur halus:
        ciri \qty{15}{\micro m}-nya sesungguhnya sebuah \emph{pita}
        selebar sekitar \qty{1}{\micro m}. Mengapa \emph{lebar} pitanya
        berperan bagi gas rumah kaca (pikirkanlah apa yang terjadi
        setelah pusat pitanya menjadi legap)?
  \item CO$_2$ yang tereksitasi secara getaran di udara jauh lebih
        mungkin kehilangan kuantumnya lewat tumbukan daripada lewat
        pemancaran (umur radiatifnya $\sim\qty{1}{s}$, waktu tumbukannya
        $\sim\qty{e-9}{s}$): ke manakah energi radiasi yang terserap itu
        sesungguhnya pergi?
  \item Sebaliknya, udara pada \qty{288}{K} mempertahankan populasi
        termal pada $n = 1$ (pertanyaan 8): apa yang dilakukan populasi
        itu yang berarti bagi neraca energinya?
\end{enumerate}

\medskip\noindent\textbf{Bagian III --- Pembukuan radiasi planetnya.}
\begin{enumerate}[resume]
  \item Matahari memancar sebagai benda \qty{5800}{K}: hitunglah puncak
        Wiennya. Apakah lenturan CO$_2$ mencegat banyak sinar matahari?
  \item Tanahnya memancar sebagai benda \qty{288}{K}: hitunglah puncak
        Wiennya, lalu tempatkanlah \qty{15}{\micro m} pada kurva termal
        itu.
  \item Jelaskanlah mekanisme rumah kacanya dalam empat kalimat: sinar
        matahari masuk, inframerah termal keluar, pencegatan pada
        \qty{15}{\micro m}, dan pemancaran ulang ke atas \emph{maupun}
        ke bawah.
  \item Pemancaran ulang yang lolos ke angkasa berasal dari lapisan
        yang tinggi dan dingin ($\sim\qty{220}{K}$): mengapa memancar
        dari lapisan yang lebih dingin mengurangi daya keluar planetnya
        pada panjang gelombang itu (ingatlah bahwa pancaran termal
        tumbuh bersama $T$)?
  \item Lebih banyak CO$_2$ mendorong lapisan pemancarnya makin tinggi
        dan makin dingin: nyatakanlah dalam satu kalimat mengapa
        permukaannya lalu harus menghangat agar pembukuannya seimbang
        kembali.
  \item Pusat pitanya sudah legap; pengaruh CO$_2$ tambahan bekerja di
        \emph{sayap} pitanya, sehingga pemaksaannya tumbuh secara
        logaritmik terhadap kepekatannya: kaitkanlah hal ini dengan
        pertanyaan 11.
\end{enumerate}

\medskip\noindent\textbf{Bagian IV --- Isotop: sidik jari
kuantumnya.}
\begin{enumerate}[resume]
  \item Frekuensi sebuah osilator berskala sebagai $\sqrt{k/\mu}$: untuk
        uluran taksetangkupnya, mengganti $^{12}$C dengan $^{13}$C
        mengubah massa efektifnya; garis teramatinya bergeser dari
        \qty{2349}{cm^{-1}} menjadi sekitar \qty{2283}{cm^{-1}}.
        Periksalah orde besar pergeseran $\approx 3\%$ ini dari massanya.
  \item Laser yang ditala pada kedua garis itu mencacah $^{13}$CO$_2$
        dan $^{12}$CO$_2$ secara terpisah dalam cuplikan gas:
        jelaskanlah mengapa tangga terkuantumkannya membuat deteksi
        terpilah isotop semacam itu mungkin sama sekali.
  \item Tumbuhan lebih menyukai isotop yang ringan, sehingga karbon
        fosil miskin $^{13}$C: apa yang telah dilakukan nisbah isotop
        terukur CO$_2$ atmosfer ketika kepekatannya naik --- dan apa
        yang dibuktikannya tentang \emph{sumber} gas tambahannya?
  \item Fisika \qty{15}{\micro m} yang sama bekerja di Venus
        ($96\%$ CO$_2$, \qty{90}{bar}): apa yang digambarkan
        permukaannya yang \qty{737}{K}?
  \item Mars juga bernapas CO$_2$ yang nyaris murni, tetapi pada
        \qty{6}{mbar}, dan ia sangat dingin: apa yang diperagakan trio
        Venus--Bumi--Mars tentang variabel mana yang mengendalikan
        kekuatan efeknya?
  \item Ringkaslah hasil bernama itu: satu kuantum sebesar \qty{83}{meV}
        --- yakni $\hbar\omega$ sebuah molekul triatom yang melentur ---
        yang terparkir pada \qty{15}{\micro m} dalam spektrum termal
        planet \qty{288}{K}, yang terserap, tertermalkan dalam hitungan
        nanosekon lalu dipancarkan ulang dari ketinggian yang dingin,
        adalah mekanisme yang membuat $0.04\%$ udara menetapkan suhu
        Bumi.
\end{enumerate}
\end{problem}
